% Part: normal-modal-logic
% Chapter: tableaux
% Section: rules-for-K

\documentclass[../../../include/open-logic-section]{subfiles}

\begin{document}

\olfileid{nml}{tab}{rul}

\olsection{\Ax{K} కోసం నియమాలు}

సాధారణ ప్రతిజ్ఞావాక్య సంయోజకాల నియమాలు, పూర్వసూచికలను
చేర్చడం మినహా, సాధారణ ప్రతిజ్ఞావాక్య చిహ్నిత
\tetoken{టాబ్లోల}{!!{tableau}s} నియమాలే. ప్రతి సందర్భంలోనూ
చిహ్నిత \tetoken{సూత్రం}{!!{formula}}~$\sFmla{S}{!A}[\sigma]$కు
నియమాన్ని వర్తింపజేస్తే, దానివల్ల వచ్చే కొత్త
\tetoken{సూత్రాలకు}{!!{formula}s} కూడా~$\sigma$ పూర్వసూచికే ఉంటుంది.
ఇది సహజంగానే అర్థమవుతుంది: ఉదాహరణకు,~$\sigma$ పేరుగల
లోకంలో $!A \land !B$ సత్యమైతే, అదే~$\sigma$లో $!A$,
$!B$ రెండూ సత్యం (వేరే ఏ లోకంలోనూ తప్పనిసరిగా కాదు).
\sourcecorrection{OLTENMLTABRUL-001}{మూలంలోని ``వేరే ఏ లోకంలోనూ కాదు''
  అనేది ఇతర లోకాలన్నింటిలో ఆ రెండు సూత్రాలు అసత్యమని చదివితే తప్పు.
  ఈ నియమం వాటిని అదే పూర్వసూచిక గల లోకంలో మాత్రమే నిష్కర్షగా
  ఇస్తుంది; ఇతర లోకాల సత్యాన్ని నిర్ణయించదు.}
ప్రతిజ్ఞావాక్య నియమాలను \olref{tab:prop-rules}లో చూపాం.

\begin{table}
  \[\def\arraystretch{3}\begin{array}{|c|c|}
    \hline
    \AxiomC{\sFmla{\True}{\lnot !A}[\sigma]}
    \RightLabel{\TRule{\True}{\lnot}}
    \UnaryInfC{\sFmla{\False}{!A}[\sigma]}
    \DisplayProof
    &
    \AxiomC{\sFmla{\False}{\lnot !A}[\sigma]}
    \RightLabel{\TRule{\False}{\lnot}}
    \UnaryInfC{\sFmla{\True}{!A}[\sigma]}
    \DisplayProof
    \\[1ex]
    \hline
    \AxiomC{\sFmla{\True}{!A \land !B}[\sigma]}
    \RightLabel{\TRule{\True}{\land}}
    \UnaryInfC{\sFmla{\True}{!A}[\sigma]}
    \noLine
    \UnaryInfC{\sFmla{\True}{!B}[\sigma]}
    \DisplayProof
    &
    \AxiomC{\sFmla{\False}{!A \land !B}[\sigma]}
    \RightLabel{\TRule{\False}{\land}}
    \UnaryInfC{$\sFmla{\False}{!A}[\sigma] \quad \mid \quad
      \sFmla{\False}{!B}[\sigma]$}
    \DisplayProof
    \\[2ex]
    \hline
    \AxiomC{\sFmla{\True}{!A \lor !B}[\sigma]}
    \RightLabel{\TRule{\True}{\lor}}
    \UnaryInfC{$\sFmla{\True}{!A}[\sigma] \quad \mid \quad
      \sFmla{\True}{!B}[\sigma]$}
    \DisplayProof
    &
    \AxiomC{\sFmla{\False}{!A \lor !B}[\sigma]}
    \RightLabel{\TRule{\False}{\lor}}
    \UnaryInfC{\sFmla{\False}{!A}[\sigma]}
    \noLine
    \UnaryInfC{\sFmla{\False}{!B}[\sigma]}
    \DisplayProof
    \\[2ex]
    \hline
    \AxiomC{\sFmla{\True}{!A \lif !B}[\sigma]}
    \RightLabel{\TRule{\True}{\lif}}
    \UnaryInfC{$\sFmla{\False}{!A}[\sigma] \quad \mid
      \quad \sFmla{\True}{!B}[\sigma]$}
    \DisplayProof
    &
    \AxiomC{\sFmla{\False}{!A \lif !B}[\sigma]}
    \RightLabel{\TRule{\False}{\lif}}
    \UnaryInfC{\sFmla{\True}{!A}[\sigma]}
    \noLine
    \UnaryInfC{\sFmla{\False}{!B}[\sigma]}
    \DisplayProof
    \\[2ex]
    \hline
  \end{array}\]
  \caption{ప్రతిజ్ఞావాక్య సంయోజకాల కోసం పూర్వసూచిక గల
    \tetoken{టాబ్లో}{!!{tableau}} నియమాలు}
  \ollabel{tab:prop-rules}
\end{table}

సంవృతత షరతు సాధారణ \tetoken{టాబ్లోల}{!!{tableau}s} మాదిరిగానే
ఉంటుంది; కానీ \tetoken{సూత్రాలు}{!!{formula}s} మాత్రమే కాక
పూర్వసూచికలు కూడా ఒకటే కావాలి. కాబట్టి ఒక శాఖలో
\[
\sFmla{\True}{!A}[\sigma] \quad\text{మరియు}\quad \sFmla{\False}{!A}[\sigma]
\]
రెండూ ఉంటే అది సంవృతమైనది; ఇక్కడ $\sigma$ ఏదో ఒక
పూర్వసూచిక,~$!A$ ఏదో ఒక \tetoken{సూత్రం}{!!{formula}}.

పరికల్పనలను ఏర్పాటు చేసే నియమాలు కూడా సాధారణ
\tetoken{టాబ్లోల}{!!{tableau}s} మాదిరిగానే ఉంటాయి; ఒక్క తేడా:
పరికల్పనలన్నింటికీ ఎప్పుడూ~$1$ పూర్వసూచికను వాడతాం.
(అన్ని పరికల్పనలకూ ఒకటే పూర్వసూచిక వాడినంతవరకు అది
ఏదైనా కావచ్చు.) కాబట్టి, ఉదాహరణకు,
\[
!B_1, \dots, !B_n \Proves !A
\]
అని అంటాం, అప్పుడు మాత్రమే కింది పరికల్పనలకు ఒక సంవృత
టాబ్లో ఉంటుంది:
\[
\sFmla{\True}{!B_1}[1], \dots, \sFmla{\True}{!B_n}[1],
\sFmla{\False}{!A}[1].
\]

\iftag{prvBox}{\iftag{prvDiamond}{మోడల్ సంయోజకాలు~$\Box$ మరియు}{మోడల్ సంయోజకం~$\Box$}}{మోడల్ సంయోజకం}\iftag{prvDiamond}{~$\Diamond$}{}
కోసం,~$\sigma$ పూర్వసూచిక గల \tetoken{ఒక సూత్రానికి}{!!a{formula}}
నియమాన్ని వర్తింపజేస్తే, దాని నిష్కర్ష పూర్వసూచిక
$\sigma.n$ అవుతుంది. అయితే ఏ~$n$ను అనుమతిస్తామనేది సంకేతం
$\True$నా, $\False$నా అన్నదానిపై ఆధారపడుతుంది.

\iftag{prvBox}{$\TRule{\True}{\Box}$ నియమం,
  $\sFmla{\True}{\Box !A}[\sigma]$ ఉన్న శాఖకు
  $\sFmla{\True}{!A}[\sigma.n]$ను చేర్చుతుంది.\iftag{prvDiamond}{ అలాగే
  }{}}{}\iftag{prvDiamond}{$\TRule{\False}{\Diamond}$ నియమం,
  $\sFmla{\False}{\Diamond !A}[\sigma]$ ఉన్న శాఖకు
  $\sFmla{\False}{!A}[\sigma.n]$ను చేర్చుతుంది.}{}
\iftag{notprvBox,notprvDiamond}{ఈ నియమాన్ని}{ఈ నియమాలను}
వర్తింపజేసే శాఖలో \emph{ఇప్పటికే} కనిపించే~$\sigma.n$
పూర్వసూచికకు మాత్రమే వర్తింపజేయవచ్చు. అటువంటి
పూర్వసూచికను (ఆ శాఖలో) ``ఉపయోగించినది'' అందాం.

\iftag{prvBox}{$\TRule{\False}{\Box}$ నియమం,
  $\sFmla{\False}{\Box !A}[\sigma]$ ఉన్న శాఖకు
  $\sFmla{\False}{!A}[\sigma.n]$ను చేర్చుతుంది.\iftag{prvDiamond}{ అలాగే
  }{}}{}\iftag{prvDiamond}{$\TRule{\True}{\Diamond}$ నియమం,
  $\sFmla{\True}{\Diamond !A}[\sigma]$ ఉన్న శాఖకు
  $\sFmla{\True}{!A}[\sigma.n]$ను చేర్చుతుంది.}{}
అయితే \iftag{notprvBox,notprvDiamond}{ఈ నియమాన్ని}{ఈ నియమాలను},
వర్తింపజేసే శాఖలో \emph{ఇప్పటివరకు కనిపించని}~$\sigma.n$
పూర్వసూచికకు మాత్రమే వర్తింపజేయవచ్చు. అటువంటి
పూర్వసూచికలను (ఆ శాఖకు) ``కొత్తవి'' అంటాం.

ఈ నియమాలు \olref{tab:rules-K}లో ఉన్నాయి.

\begin{table}
  \begin{center}
    \def\arraystretch{3}\def\fCenter{}
    \begin{tabular}{|c|c|}
    \hline
    \iftag{prvBox}{
      \AxiomC{\sFmla{\True}{\Box !A}[\sigma]}
      \RightLabel{\TRule{\True}{\Box}}
      \UnaryInfC{\sFmla{\True}{!A}[\sigma.n]}
      \DisplayProof
      &
      \AxiomC{\sFmla{\False}{\Box !A}[\sigma]}
      \RightLabel{\TRule{\False}{\Box}}
      \UnaryInfC{\sFmla{\False}{!A}[\sigma.n]}
      \DisplayProof\\
      $\sigma.n$ ఉపయోగించినది & $\sigma.n$ కొత్తది
      \\[1ex]
      \hline}{}
    \iftag{prvDiamond}{
      \AxiomC{\sFmla{\True}{\Diamond !A}[\sigma]}
      \RightLabel{\TRule{\True}{\Diamond}}
      \UnaryInfC{\sFmla{\True}{!A}[\sigma.n]}
      \DisplayProof
      &
      \AxiomC{\sFmla{\False}{\Diamond !A}[\sigma]}
      \RightLabel{\TRule{\False}{\Diamond}}
      \UnaryInfC{\sFmla{\False}{!A}[\sigma.n]}
      \DisplayProof\\
      $\sigma.n$ కొత్తది & $\sigma.n$ ఉపయోగించినది
      \\[1ex]
      \hline}{}
    \end{tabular}
  \end{center}
  \caption{\Ax{K} కోసం మోడల్ నియమాలు.}
  \ollabel{tab:rules-K}
\end{table}

\iftag{prvBox}{\TRule{\True}{\Box}}{\TRule{\False}{\Diamond}}
నియమానికి ఉపయోగించిన పూర్వసూచికే ఉండాలనే ఆంక్ష అవసరం;
లేకపోతే కిందిదాన్ని సంవృత \tetoken{టాబ్లో}{!!{tableau}}గా
లెక్కించాల్సి వస్తుంది:
\iftag{notprvBox,notprvDiamond}{%
  \iftag{prvBox}{%
  \begin{oltableau}
    [\pFmla{\True}{\Box \formula{A}}{1}, just = \TAss
      [\pFmla{\False}{\lnot\Box\lnot \formula{A}}{1}, just = \TAss
        [\pFmla{\True}{\formula{A}}{1.1}, just = {\TRule{\True}{\Box}[1]}
          [\pFmla{\True}{\Box\lnot\formula{A}}{1},
            just ={\TRule{\False}{\lnot}[2]}
            [\pFmla{\True}{\lnot\formula{A}}{1.1},
              just = {\TRule{\True}{\Box}[4]}
              [\pFmla{\False}{\formula{A}}{1.1},
                  just ={\TRule{\True}{\lnot}[5]}, close]
            ]
          ]
        ]
      ]
    ]
  \end{oltableau}
  }{
  \begin{oltableau}
    [\pFmla{\True}{\lnot\Diamond\lnot \formula{A}}{1}, just = \TAss
      [\pFmla{\False}{\Diamond\formula{A}}{1}, just = \TAss
        [\pFmla{\False}{\formula{A}}{1.1}, just = {\TRule{\False}{\Diamond}[2]}
          [\pFmla{\False}{\Diamond\lnot\formula{A}}{1},
            just ={\TRule{\True}{\lnot}[1]}
            [\pFmla{\False}{\lnot\formula{A}}{1.1},
              just = {\TRule{\False}{\Diamond}[4]}
              [\pFmla{\True}{\formula{A}}{1.1},
                  just ={\TRule{\True}{\lnot}[5]}, close]
            ]
          ]
        ]
      ]
    ]
  \end{oltableau}
}}{
  \begin{oltableau}
    [\pFmla{\True}{\Box \formula{A}}{1}, just = \TAss
      [\pFmla{\False}{\Diamond \formula{A}}{1}, just = \TAss
        [\pFmla{\True}{\formula{A}}{1.1}, just = {\TRule{\True}{\Box}[1]}
          [\pFmla{\False}{\formula{A}}{1.1},
            just = {\TRule{\False}{\Diamond}[2]}, close]
        ]
      ]
    ]
\end{oltableau}}

కానీ $\Box \formula{A} \Entails/ \Diamond \formula{A}$; కాబట్టి
అలా అయితే మన నిరూపణ వ్యవస్థ నిర్దుష్టం కాదు. అలాగే,
$\Diamond \formula{A} \Entails/ \Box \formula{A}$; అయితే
\iftag{prvBox}{\TRule{\False}{\Box}}{\TRule{\True}{\Diamond}}
నియమం కోసం పూర్వసూచిక కొత్తదిగా ఉండాలనే ఆంక్ష లేకపోతే,
కిందిది సంవృత టాబ్లో అవుతుంది:
\iftag{notprvBox,notprvDiamond}{%
  \iftag{prvBox}{%
  \sourcecorrection{OLTENMLTABRUL-002}{స్థిర మూలంలోని కేవలం
    Box ఉన్న శాఖ చిత్రం, \sFmla{\False}{\Box !A}[1] నుంచి
    \sFmla{\False}{!A}[1.1]కి వెళ్లే నియమాన్ని పొరపాటుగా
    \TRule{\True}{\Box} అని గుర్తించింది. దాన్ని
    \TRule{\False}{\Box}గా సరిచేశాం; కొత్త పూర్వసూచిక
    షరతు ఉల్లంఘనను చూపే మిగతా శాఖ యథాతథం.}
  \begin{oltableau}
    [\pFmla{\True}{\lnot\Box\lnot \formula{A}}{1}, just = \TAss
      [\pFmla{\False}{\Box \formula{A}}{1}, just = \TAss
        [\pFmla{\False}{\formula{A}}{1.1}, just = {\TRule{\False}{\Box}[2]}
          [\pFmla{\False}{\Box\lnot\formula{A}}{1},
            just ={\TRule{\True}{\lnot}[1]}
            [\pFmla{\False}{\lnot\formula{A}}{1.1},
              just = {\TRule{\False}{\Box}[4]}
              [\pFmla{\True}{\formula{A}}{1.1},
                  just ={\TRule{\False}{\lnot}[5]}, close]
            ]
          ]
        ]
      ]
    ]
  \end{oltableau}
  }{
  \begin{oltableau}
    [\pFmla{\True}{\Diamond \formula{A}}{1}, just = \TAss
      [\pFmla{\False}{\lnot\Diamond\lnot\formula{A}}{1}, just = \TAss
        [\pFmla{\True}{\formula{A}}{1.1}, just = {\TRule{\True}{\Diamond}[1]}
          [\pFmla{\True}{\Diamond\lnot\formula{A}}{1},
            just ={\TRule{\False}{\lnot}[2]}
            [\pFmla{\True}{\lnot\formula{A}}{1.1},
              just = {\TRule{\True}{\Diamond}[4]}
              [\pFmla{\False}{\formula{A}}{1.1},
                  just ={\TRule{\True}{\lnot}[5]}, close]
            ]
          ]
        ]
      ]
    ]
  \end{oltableau}
}}{
  \begin{oltableau}
    [\pFmla{\True}{\Diamond \formula{A}}{1}, just = \TAss,name=one
      [\pFmla{\False}{\Box \formula{A}}{1}, just = \TAss, name=two
        [\pFmla{\True}{\formula{A}}{1.1}, just = {\TRule{\True}{\Diamond}[1]}
          [\pFmla{\False}{\formula{A}}{1.1},
            just = {\TRule{\False}{\Box}[2]}, close]
        ]
      ]
    ]
  \end{oltableau}}

\end{document}
